
@misc{carrisi_structured_2025,
	title = {A {Structured} {Unplugged} {Approach} for {Foundational} {AI} {Literacy} in {Primary} {Education}},
	url = {http://arxiv.org/abs/2505.21398},
	doi = {10.48550/arXiv.2505.21398},
	abstract = {Younger generations are growing up in a world increasingly shaped by intelligent technologies, making early 
AI literacy crucial for developing the skills to critically understand and navigate them. However, education in this field often 
emphasizes tool-based learning, prioritizing usage over understanding the underlying concepts. This lack of knowledge leaves 
non-experts, especially children, prone to misconceptions, unrealistic expectations, and difficulties in recognizing biases and 
stereotypes. In this paper, we propose a structured and replicable teaching approach that fosters foundational AI literacy in 
primary students, by building upon core mathematical elements closely connected to and of interest in primary curricula, to 
strengthen conceptualization, data representation, classification reasoning, and evaluation of AI. To assess the effectiveness 
of our approach, we conducted an empirical study with thirty-one fifth-grade students across two classes, evaluating their 
progress through a post-test and a satisfaction survey. Our results indicate improvements in terminology understanding and 
usage, features description, logical reasoning, and evaluative skills, with students showing a deeper comprehension of 
decision-making processes and their limitations. Moreover, the approach proved engaging, with students particularly enjoying 
activities that linked AI concepts to real-world reasoning. Materials: https://github.com/tail-unica/ai-literacy-primary-ed.},
	urldate = {2025-10-07},
	publisher = {arXiv},
	author = {Carrisi, Maria Cristina and Marras, Mirko and Vergallo, Sara},
	month = may,
	year = {2025},
	note = {arXiv:2505.21398 [cs]},
	keywords = {Computer Science - Artificial Intelligence, Computer Science - Emerging Technologies},
	annote = {Comment: Under review},
	annote = {Comment: Under review},
	file = {Preprint PDF:C\:\\Users\\cbenahmed\\Zotero\\storage\\FDXR4PVJ\\Carrisi et al. - 2025 - A Structured Unplugged 
Approach for Foundational AI Literacy in Primary 
Education.pdf:application/pdf;Snapshot:C\:\\Users\\cbenahmed\\Zotero\\storage\\9DSVGGRU\\2505.html:text/html},
}

@inproceedings{dybboe_bit_2025,
	title = {Bit: sort : {Bringing} {Tangible} {Computing} to {Computer} {Science} {Unplugged} to {Support} {Children}'s 
{Algorithmic} {Explorations}},
	shorttitle = {Bit},
	doi = {10.1145/3713043.3728854},
	booktitle = {Proceedings of the 24th {Interaction} {Design} and {Children}, {IDC} 2025, {Reykjavik}, {Iceland}, {June} 
23-26, 2025},
	publisher = {ACM},
	author = {Dybboe, Maja and Kaspersen, Magnus Høholt and Bjerrum, Johanne Birkkjær and Petersen, Marianne Graves},
	year = {2025},
	keywords = {Article récent},
	pages = {429--443},
}

@article{wada_teaching_2024,
	title = {Teaching {Computer} {Science} {Unplugged} in {Online} for {Undergraduate} {College} {Students}},
	volume = {32},
	url = {https://doi.org/10.2197/ipsjjip.32.127},
	doi = {10.2197/IPSJJIP.32.127},
	urldate = {2025-10-07},
	journal = {J. Inf. Process.},
	author = {Wada, Ben Tsutom},
	year = {2024},
	pages = {127--132},
	file = {Texte intégral:C\:\\Users\\cbenahmed\\Zotero\\storage\\598TNS4B\\Wada - 2024 - Teaching Computer Science 
Unplugged in Online for Undergraduate College Students.pdf:application/pdf},
}

@inproceedings{jinnouchi_visualizing_2023,
	title = {Visualizing {Program} {Behavior} with a {Ball} and {Pipes} for {Computer} {Science} {Unplugged}},
	doi = {10.1109/APSEC60848.2023.00099},
	booktitle = {30th {Asia}-{Pacific} {Software} {Engineering} {Conference}, {APSEC} 2023, {Seoul}, {Republic} of {Korea}, 
{December} 4-7, 2023},
	publisher = {IEEE},
	author = {Jinnouchi, Snmika and Tsunoda, Masateru},
	year = {2023},
	pages = {663--664},
}

@inproceedings{landman_reshaping_2023,
	series = {Lecture {Notes} in {Computer} {Science}},
	title = {Reshaping {Unplugged} {Computer} {Science} {Workshops} for {Primary} {School} {Education}},
	volume = {14296},
	url = {https://doi.org/10.1007/978-3-031-44900-0\_11},
	doi = {10.1007/978-3-031-44900-0_11},
	urldate = {2025-10-07},
	booktitle = {Informatics in {Schools}. {Beyond} {Bits} and {Bytes}: {Nurturing} {Informatics} {Intelligence} in 
{Education} - 16th {International} {Conference} on {Informatics} in {Schools}: {Situation}, {Evolution}, and {Perspectives}, 
{ISSEP} 2023, {Lausanne}, {Switzerland}, {October} 23-25, 2023, {Proceedings}},
	publisher = {Springer},
	author = {Landman, Martina and Rain, Sophie and Kovács, Laura and Futschek, Gerald},
	editor = {Pellet, Jean-Philippe and Parriaux, Gabriel},
	year = {2023},
	pages = {139--151},
}

@article{huang_critical_2021,
	title = {A critical review of literature on "unplugged" pedagogies in {K}-12 computer science and computational thinking 
education},
	volume = {31},
	doi = {10.1080/08993408.2020.1789411},
	number = {1},
	journal = {Comput. Sci. Educ.},
	author = {Huang, Wendy and Looi, Chee-Kit},
	year = {2021},
	pages = {83--111},
}

@article{yildiz_computer_2021,
	title = {A {Computer} {Science} {Unplugged} {Activity}: {CityMap}},
	volume = {5},
	shorttitle = {A {Computer} {Science} {Unplugged} {Activity}},
	url = {https://doi.org/10.21585/ijcses.v5i2.110},
	doi = {10.21585/IJCSES.V5I2.110},
	number = {2},
	urldate = {2025-10-07},
	journal = {Int. J. Comput. Sci. Educ. Sch.},
	author = {Yildiz, Merve and Karal, Hasan},
	year = {2021},
	pages = {14--27},
	file = {Texte intégral:C\:\\Users\\cbenahmed\\Zotero\\storage\\VQU6WMV7\\Yildiz et Karal - 2021 - A Computer Science 
Unplugged Activity CityMap.pdf:application/pdf},
}

@inproceedings{svabensky_stack_2021,
	title = {The {Stack}: {Unplugged} {Activities} for {Teaching} {Computer} {Science}},
	shorttitle = {The {Stack}},
	doi = {10.1145/3408877.3439569},
	booktitle = {{SIGCSE} '21: {The} 52nd {ACM} {Technical} {Symposium} on {Computer} {Science} {Education}, {Virtual} 
{Event}, {USA}, {March} 13-20, 2021},
	publisher = {ACM},
	author = {Svábenský, Valdemar and Ukrop, Martin},
	editor = {Sherriff, Mark and Merkle, Laurence D. and Cutter, Pamela A. and Monge, Alvaro E. and Sheard, Judithe},
	year = {2021},
	pages = {1242},
	file = {Texte intégral:C\:\\Users\\cbenahmed\\Zotero\\storage\\NZH8SE64\\Svábenský et Ukrop - 2021 - The Stack Unplugged 
Activities for Teaching Computer Science.pdf:application/pdf},
}

@inproceedings{bell_teaching_2020,
	title = {Teaching {Teachers} to {Teach} {Computer} {Science} - {Unplugged} or {Plugged}-in?},
	doi = {10.1145/3372782.3406284},
	booktitle = {{ICER} 2020: {International} {Computing} {Education} {Research} {Conference}, {Virtual} {Event}, {New} 
{Zealand}, {August} 10-12, 2020},
	publisher = {ACM},
	author = {Bell, Tim},
	editor = {Robins, Anthony V. and Moskal, Adon and Ko, Amy J. and McCauley, Renée},
	year = {2020},
	pages = {1},
}

%Ref 4
@article{sendurur_investigation_2019,
	title = {Investigation of pre-service computer science Teachers' CS-unplugged design practices},
	volume = {24},
	number = {6},
	author = {P. Sendurur},
	pages = {3823--3840},
}

%Ref 6
@inproceedings{srivastava_assessing_2019,
	title = {Assessing the Integration of Parallel and Distributed Computing in Early Undergraduate Computer Science 
Curriculum using Unplugged Activities},
	booktitle = {2019 IEEE/ACM Workshop on Education for High-Performance Computing, EduHPC@SC 2019, Denver, CO, USA, 
November 17, 2019},
	publisher = {IEEE},
	author = {Srivastava, Srishti and Smith, Mary L. and Ghimire, Amrita and Gao, Shaun},
	year = {2019},
	pages = {17--24},
}

@article{fees_unplugged_2018,
	title = {Unplugged cybersecurity: {An} approach for bringing computer science into the classroom},
	volume = {2},
	shorttitle = {Unplugged cybersecurity},
	url = {https://doi.org/10.21585/ijcses.v2i1.21},
	doi = {10.21585/IJCSES.V2I1.21},
	number = {1},
	urldate = {2025-10-07},
	journal = {Int. J. Comput. Sci. Educ. Sch.},
	author = {Fees, Rachel E. and Rosa, Jennifer A. Da and Durkin, Sarah S. and Murray, Mark M. and Moran, Angela L.},
	year = {2018},
	keywords = {Article avant 2020, Article récent},
	pages = {3--13},
}

@inproceedings{thies_back_2016,
	title = {Back to {School}: {Computer} {Science} {Unplugged} in the {Wild}},
	shorttitle = {Back to {School}},
	doi = {10.1145/2899415.2899442},
	booktitle = {Proceedings of the 2016 {ACM} {Conference} on {Innovation} and {Technology} in {Computer} {Science} 
{Education}, {ITiCSE} 2016, {Arequipa}, {Peru}, {July} 9-13, 2016},
	publisher = {ACM},
	author = {Thies, Renate and Vahrenhold, Jan},
	editor = {Clear, Alison and Cuadros-Vargas, Ernesto and Carter, Janet and Túpac, Yván},
	year = {2016},
	pages = {118--123},
}

@inproceedings{wohl_teaching_2015,
	title = {Teaching {Computer} {Science} to 5-7 year-olds: {An} initial study with {Scratch}, {Cubelets} and unplugged 
computing},
	shorttitle = {Teaching {Computer} {Science} to 5-7 year-olds},
	doi = {10.1145/2818314.2818340},
	booktitle = {Proceedings of the {Workshop} in {Primary} and {Secondary} {Computing} {Education}, {WiPSCE} 2015, 
{London}, {United} {Kingdom}, {November} 9-11, 2015},
	publisher = {ACM},
	author = {Wohl, Benjamin S. and Porter, Barry and Clinch, Sarah},
	editor = {Gal-Ezer, Judith and Sentance, Sue and Vahrenhold, Jan},
	year = {2015},
	pages = {55--60},
	file = {Version acceptée:C\:\\Users\\cbenahmed\\Zotero\\storage\\Z5HJFKP4\\Wohl et al. - 2015 - Teaching Computer 
Science to 5-7 year-olds An initial study with Scratch, Cubelets and unplugged co.pdf:application/pdf},
}

@inproceedings{marghitu_kodu_2014,
	title = {Kodu alice and computer science unplugged: a model of effective introducing middle school students to computer 
science and computational thinking (abstract only)},
	shorttitle = {Kodu alice and computer science unplugged},
	doi = {10.1145/2538862.2544300},
	booktitle = {The 45th {ACM} {Technical} {Symposium} on {Computer} {Science} {Education}, {SIGCSE} 2014, {Atlanta}, {GA}, 
{USA}, {March} 5-8, 2014},
	publisher = {ACM},
	author = {Marghitu, Daniela and Thomas, Lavaris and Rawajfih, Yasmeen and Hall, Jillian and Marshall, Andrew},
	editor = {Dougherty, J. D. and Nagel, Kris and Decker, Adrienne and Eiselt, Kurt},
	year = {2014},
	pages = {712},
}

%Ref 2
@inproceedings{bell_computer_2012,
	series = {Lecture Notes in Computer Science},
	title = {Computer Science Unplugged and {Related} Projects in Math and Computer Science Popularization.},
	volume = {7370},
	booktitle = {The Multivariate Algorithmic Revolution and Beyond - Essays Dedicated to Michael R. Fellows on the Occasion 
of His 60th Birthday},
	publisher = {Springer},
	author = {T.Bell, F. A. Rosamond, N.Casey},
	editor = {Bodlaender, Hans L. and Downey, Rod and Fomin, Fedor V. and Marx, Dániel},
	pages = {398--456},
}

@inproceedings{marghitu_teaching_2012,
	title = {Teaching computer science and programming concepts using {LEGO} {NXT} and {TETRIX} robotics, and computer 
science unplugged activities (abstract only)},
	doi = {10.1145/2157136.2157394},
	booktitle = {Proceedings of the 43rd {ACM} technical symposium on {Computer} science education, {SIGCSE} 2012, 
{Raleigh}, {NC}, {USA}, {February} 29 - {March} 3, 2012},
	publisher = {ACM},
	author = {Marghitu, Daniela and Brahim, Taha Ben and Weaver, John},
	editor = {King, Laurie A. Smith and Musicant, David R. and Camp, Tracy and Tymann, Paul T.},
	year = {2012},
	pages = {671},
}

@article{cortina_computer_2009,
	title = {Computer science unplugged: pre-conference workshop},
	volume = {24},
	shorttitle = {Computer science unplugged},
	url = {https://dl.acm.org/doi/10.5555/1516595.1516618},
	doi = {10.5555/1516595.1516618},
	number = {5},
	urldate = {2025-10-07},
	journal = {J. Comput. Sci. Coll.},
	author = {Cortina, Thomas J.},
	year = {2009},
	pages = {107},
}

@article{henderson_computer_2009,
	title = {Computer science unplugged: tutorial presentation},
	volume = {24},
	shorttitle = {Computer science unplugged},
	url = {https://dl.acm.org/doi/10.5555/1409873.1409903},
	doi = {10.5555/1409873.1409903},
	number = {3},
	urldate = {2025-10-07},
	journal = {J. Comput. Sci. Coll.},
	author = {Henderson, Peter B. and Lambert, Lynn},
	year = {2009},
	pages = {159},
}

@article{lucas_k-6_2009,
	title = {K-6 outreach using "{Computer} {Science} {Unplugged}"},
	volume = {24},
	url = {https://dl.acm.org/doi/10.5555/1529995.1530008},
	doi = {10.5555/1529995.1530008},
	number = {6},
	urldate = {2025-10-07},
	journal = {J. Comput. Sci. Coll.},
	author = {Lucas, Joan M.},
	year = {2009},
	pages = {62--63},
}

@article{henderson_computer_2008,
	title = {Computer {Science} unplugged},
	volume = {23},
	url = {https://dl.acm.org/doi/10.5555/1295109.1295146},
	doi = {10.5555/1295109.1295146},
	number = {3},
	urldate = {2025-10-07},
	journal = {J. Comput. Sci. Coll.},
	author = {Henderson, Peter B.},
	year = {2008},
	pages = {168},
}

%Ref 3
@inproceedings{nishida_new_2008,
	series = {Lecture Notes in Computer Science},
	title = {New Methodology of Information Education with "Computer Science Unplugged"},
	volume = {5090},
	booktitle = {Informatics {Education} - {Supporting} {Computational} {Thinking}, {Third} {International} {Conference} on 
{Informatics} in {Secondary} {Schools} - {Evolution} and {Perspectives}, {ISSEP} 2008, {Torun}, {Poland}, {July} 1-4, 2008, 
{Proceedings}},
	publisher = {Springer},
	author = {Nishida, Tomohiro and Idosaka, Yukio and Hofuku, Yayoi and Kanemune, Susumu and Kuno, Yasushi},
	editor = {Mittermeir, Roland T. and Syslo, Maciej M.},
	year = {2008},
	pages = {241--252},
}

@book{bell_computer_2002,
	title = {Computer {Science} {Unplugged} - an enrichment and extension programme for primary-aged children},
	url = {http://csunplugged.org/},
	urldate = {2025-10-07},
	publisher = {csunplugged.org},
	author = {Bell, Tim and Witten, Ian H. and Fellows, Michael R.},
	year = {2002},
}

@misc{morales-navarro_what_2025,
	title = {What {Can} {Youth} {Learn} {About} {Artificial} {Intelligence} and {Machine} {Learning} in {One} {Hour}? 
{Examining} {How} {Hour} of {Code} {Activities} {Address} the {Five} {Big} {Ideas} of {AI}},
	shorttitle = {What {Can} {Youth} {Learn} {About} {Artificial} {Intelligence} and {Machine} {Learning} in {One} {Hour}?},
	url = {http://arxiv.org/abs/2412.11911},
	doi = {10.48550/arXiv.2412.11911},
	abstract = {The prominence of artificial intelligence and machine learning in everyday life has led to efforts to foster 
AI literacy for all K-12 students. In this paper, we review how Hour of Code activities engage with the five big ideas of AI, in 
particular with machine learning and societal impact. We found that a large majority of activities focus on perception and 
machine learning, with little attention paid to representation and other topics. A surprising finding was the increased 
attention paid to critical aspects of computing. However, we also observed a limited engagement with hands-on activities. In the 
discussion, we address how future introductory activities could be designed to offer a broader array of topics, including the 
development of tools to introduce novices to artificial intelligence and machine learning and the design of more unplugged and 
collaborative activities.},
	urldate = {2025-10-07},
	publisher = {arXiv},
	author = {Morales-Navarro, Luis and Kafai, Yasmin B. and Yang, Eric and Suryana, Asep},
	month = jan,
	year = {2025},
	note = {arXiv:2412.11911 [cs]},
	keywords = {Computer Science - Computers and Society},
	file = {Preprint PDF:C\:\\Users\\cbenahmed\\Zotero\\storage\\FRWIARES\\Morales-Navarro et al. - 2025 - What Can Youth 
Learn About Artificial Intelligence and Machine Learning in One Hour Examining How H.pdf:application/pdf},
}

@misc{jinnouchi_influence_2024,
	title = {Influence of {Skill} and {Knowledge} of {Programmers} on {Program} {Behavior} {Visualization} by {CS} 
{Unplugged}},
	url = {http://arxiv.org/abs/2412.01831},
	doi = {10.48550/arXiv.2412.01831},
	abstract = {Computer science unplugged (CS unplugged) is a method of teaching computer science and computational 
thinking. It does not use a computer but employs physical materials. As CS unplugged, past studies proposed a new method to 
visualize programming behavior, and evaluated the method based on the understanding test of program. However, the past studied 
did not consider the factors affecting the test such as participants knowledge. This study analyzed the relationship between the 
factors and the test results.},
	urldate = {2025-10-07},
	publisher = {arXiv},
	author = {Jinnouchi, Sumika and Tsunoda, Masateru},
	month = nov,
	year = {2024},
	note = {arXiv:2412.01831 [cs]},
	keywords = {Computer Science - Computers and Society, Computer Science - Software Engineering},
	annote = {Comment: 2 pages, 3 figures, 2 tables},
	file = {Preprint PDF:C\:\\Users\\cbenahmed\\Zotero\\storage\\GLJUHQNJ\\Jinnouchi et Tsunoda - 2024 - Influence of Skill 
and Knowledge of Programmers on Program Behavior Visualization by CS Unplugged.pdf:application/pdf},
}

@misc{mathieu_alphastar_2023,
	title = {{AlphaStar} {Unplugged}: {Large}-{Scale} {Offline} {Reinforcement} {Learning}},
	shorttitle = {{AlphaStar} {Unplugged}},
	url = {http://arxiv.org/abs/2308.03526},
	doi = {10.48550/arXiv.2308.03526},
	abstract = {StarCraft II is one of the most challenging simulated reinforcement learning environments; it is partially 
observable, stochastic, multi-agent, and mastering StarCraft II requires strategic planning over long time horizons with 
real-time low-level execution. It also has an active professional competitive scene. StarCraft II is uniquely suited for 
advancing offline RL algorithms, both because of its challenging nature and because Blizzard has released a massive dataset of 
millions of StarCraft II games played by human players. This paper leverages that and establishes a benchmark, called AlphaStar 
Unplugged, introducing unprecedented challenges for offline reinforcement learning. We define a dataset (a subset of Blizzard's 
release), tools standardizing an API for machine learning methods, and an evaluation protocol. We also present baseline agents, 
including behavior cloning, offline variants of actor-critic and MuZero. We improve the state of the art of agents using only 
offline data, and we achieve 90\% win rate against previously published AlphaStar behavior cloning agent.},
	urldate = {2025-10-07},
	publisher = {arXiv},
	author = {Mathieu, Michaël and Ozair, Sherjil and Srinivasan, Srivatsan and Gulcehre, Caglar and Zhang, Shangtong and 
Jiang, Ray and Paine, Tom Le and Powell, Richard and Żołna, Konrad and Schrittwieser, Julian and Choi, David and Georgiev, Petko 
and Toyama, Daniel and Huang, Aja and Ring, Roman and Babuschkin, Igor and Ewalds, Timo and Bordbar, Mahyar and Henderson, Sarah 
and Colmenarejo, Sergio Gómez and Oord, Aäron van den and Czarnecki, Wojciech Marian and Freitas, Nando de and Vinyals, Oriol},
	month = aug,
	year = {2023},
	note = {arXiv:2308.03526 [cs]},
	keywords = {Computer Science - Artificial Intelligence, Computer Science - Machine Learning},
	annote = {Comment: 32 pages, 13 figures, previous version published as a NeurIPS 2021 workshop: 
https://openreview.net/forum?id=Np8Pumfoty},
	file = {Preprint PDF:C\:\\Users\\cbenahmed\\Zotero\\storage\\YRMBXA62\\Mathieu et al. - 2023 - AlphaStar Unplugged 
Large-Scale Offline Reinforcement Learning.pdf:application/pdf},
}

@misc{michaeli_data_2023,
	title = {Data, {Trees}, and {Forests} -- {Decision} {Tree} {Learning} in {K}-12 {Education}},
	url = {http://arxiv.org/abs/2305.06442},
	doi = {10.48550/arXiv.2305.06442},
	abstract = {As a consequence of the increasing influence of machine learning on our lives, everyone needs competencies 
to understand corresponding phenomena, but also to get involved in shaping our world and making informed decisions regarding the 
influences on our society. Therefore, in K-12 education, students need to learn about core ideas and principles of machine 
learning. However, for this target group, achieving all of the aforementioned goals presents an enormous challenge. To this end, 
we present a teaching concept that combines a playful and accessible unplugged approach focusing on conceptual understanding 
with empowering students to actively apply machine learning methods and reflect their influence on society, building upon 
decision tree learning.},
	urldate = {2025-10-07},
	publisher = {arXiv},
	author = {Michaeli, Tilman and Seegerer, Stefan and Kerber, Lennard and Romeike, Ralf},
	month = may,
	year = {2023},
	note = {arXiv:2305.06442 [cs]},
	keywords = {Computer Science - Computers and Society, Computer Science - Machine Learning},
	annote = {Comment: To be published in the proceedings of the 3rd Teaching in Machine Learning Workshop, PMLR, 2022},
	file = {Preprint PDF:C\:\\Users\\cbenahmed\\Zotero\\storage\\8Z79JNMJ\\Michaeli et al. - 2023 - Data, Trees, and Forests 
-- Decision Tree Learning in K-12 Education.pdf:application/pdf},
}

@misc{liu_efficient_2023,
	title = {Efficient {Offline} {Policy} {Optimization} with a {Learned} {Model}},
	url = {http://arxiv.org/abs/2210.05980},
	doi = {10.48550/arXiv.2210.05980},
	abstract = {MuZero Unplugged presents a promising approach for offline policy learning from logged data. It conducts 
Monte-Carlo Tree Search (MCTS) with a learned model and leverages Reanalyze algorithm to learn purely from offline data. For 
good performance, MCTS requires accurate learned models and a large number of simulations, thus costing huge computing time. 
This paper investigates a few hypotheses where MuZero Unplugged may not work well under the offline RL settings, including 1) 
learning with limited data coverage; 2) learning from offline data of stochastic environments; 3) improperly parameterized 
models given the offline data; 4) with a low compute budget. We propose to use a regularized one-step look-ahead approach to 
tackle the above issues. Instead of planning with the expensive MCTS, we use the learned model to construct an advantage 
estimation based on a one-step rollout. Policy improvements are towards the direction that maximizes the estimated advantage 
with regularization of the dataset. We conduct extensive empirical studies with BSuite environments to verify the hypotheses and 
then run our algorithm on the RL Unplugged Atari benchmark. Experimental results show that our proposed approach achieves stable 
performance even with an inaccurate learned model. On the large-scale Atari benchmark, the proposed method outperforms MuZero 
Unplugged by 43\%. Most significantly, it uses only 5.6\% wall-clock time (i.e., 1 hour) compared to MuZero Unplugged (i.e., 
17.8 hours) to achieve a 150\% IQM normalized score with the same hardware and software stacks. Our implementation is 
open-sourced at https://github.com/sail-sg/rosmo.},
	urldate = {2025-10-07},
	publisher = {arXiv},
	author = {Liu, Zichen and Li, Siyi and Lee, Wee Sun and Yan, Shuicheng and Xu, Zhongwen},
	month = feb,
	year = {2023},
	note = {arXiv:2210.05980 [cs]},
	keywords = {Computer Science - Machine Learning},
	annote = {Comment: ICLR2023},
	file = {Preprint PDF:C\:\\Users\\cbenahmed\\Zotero\\storage\\9HEGWYAV\\Liu et al. - 2023 - Efficient Offline Policy 
Optimization with a Learned Model.pdf:application/pdf},
}

@misc{ghasemipour_why_2022,
	title = {Why {So} {Pessimistic}? {Estimating} {Uncertainties} for {Offline} {RL} through {Ensembles}, and {Why} {Their} 
{Independence} {Matters}},
	shorttitle = {Why {So} {Pessimistic}?},
	url = {http://arxiv.org/abs/2205.13703},
	doi = {10.48550/arXiv.2205.13703},
	abstract = {Motivated by the success of ensembles for uncertainty estimation in supervised learning, we take a renewed 
look at how ensembles of \$Q\$-functions can be leveraged as the primary source of pessimism for offline reinforcement learning 
(RL). We begin by identifying a critical flaw in a popular algorithmic choice used by many ensemble-based RL algorithms, namely 
the use of shared pessimistic target values when computing each ensemble member's Bellman error. Through theoretical analyses 
and construction of examples in toy MDPs, we demonstrate that shared pessimistic targets can paradoxically lead to value 
estimates that are effectively optimistic. Given this result, we propose MSG, a practical offline RL algorithm that trains an 
ensemble of \$Q\$-functions with independently computed targets based on completely separate networks, and optimizes a policy 
with respect to the lower confidence bound of predicted action values. Our experiments on the popular D4RL and RL Unplugged 
offline RL benchmarks demonstrate that on challenging domains such as antmazes, MSG with deep ensembles surpasses highly 
well-tuned state-of-the-art methods by a wide margin. Additionally, through ablations on benchmarks domains, we verify the 
critical significance of using independently trained \$Q\$-functions, and study the role of ensemble size. Finally, as using 
separate networks per ensemble member can become computationally costly with larger neural network architectures, we investigate 
whether efficient ensemble approximations developed for supervised learning can be similarly effective, and demonstrate that 
they do not match the performance and robustness of MSG with separate networks, highlighting the need for new efforts into 
efficient uncertainty estimation directed at RL.},
	urldate = {2025-10-07},
	publisher = {arXiv},
	author = {Ghasemipour, Seyed Kamyar Seyed and Gu, Shixiang Shane and Nachum, Ofir},
	month = may,
	year = {2022},
	note = {arXiv:2205.13703 [cs]},
	keywords = {Computer Science - Machine Learning},
	annote = {Comment: Our codebase can be found at https://github.com/google-research/google-research/tree/master/jrl},
	file = {Preprint PDF:C\:\\Users\\cbenahmed\\Zotero\\storage\\IRXF4QSP\\Ghasemipour et al. - 2022 - Why So Pessimistic 
Estimating Uncertainties for Offline RL through Ensembles, and Why Their Indepen.pdf:application/pdf},
}

@misc{capecchi_visual_2022,
	title = {Visual and unplugged coding lessons with smart toys},
	url = {http://arxiv.org/abs/2205.11644},
	doi = {10.48550/arXiv.2205.11644},
	abstract = {Our Computer science k-12 education research group and the educational toy company Quercetti have been 
collaborating together to design and manufacture toys that help stimulate and consolidate so-called computational thinking. This 
approach is inspired by methods already consolidated in the literature and widespread worldwide such as the Bebras tasks and 
CS-Unplugged. This paper describes two smart toys, their design process, educational activities that can be proposed by teachers 
exploiting the two toys, the evaluation's results from some teachers, and finally feedback and reviews from buyers. The main 
activities proposed by these toys leverage visual coding through small colored physical items (e.g., pegs and balls) to deliver 
the unplugged activities to young users.},
	urldate = {2025-10-07},
	publisher = {arXiv},
	author = {Capecchi, Sara and Gena, Cristina and Lombardi, Ilaria},
	month = may,
	year = {2022},
	note = {arXiv:2205.11644 [cs]},
	keywords = {Computer Science - Computers and Society},
	annote = {Comment: Extented version of the short paper accepted at AVI 2022 (Visual and unplugged coding with smart 
toys)},
	file = {Preprint PDF:C\:\\Users\\cbenahmed\\Zotero\\storage\\2AAIH2Y6\\Capecchi et al. - 2022 - Visual and unplugged 
coding lessons with smart toys.pdf:application/pdf},
}

@article{el-hamamsy_competent_2022,
	title = {The competent {Computational} {Thinking} test ({cCTt}): {Development} and validation of an unplugged 
{Computational} {Thinking} test for upper primary school},
	volume = {60},
	issn = {0735-6331, 1541-4140},
	shorttitle = {The competent {Computational} {Thinking} test ({cCTt})},
	url = {http://arxiv.org/abs/2203.05980},
	doi = {10.1177/07356331221081753},
	abstract = {With the increasing importance of Computational Thinking (CT) at all levels of education, it is essential to 
have valid and reliable assessments. Currently, there is a lack of such assessments in upper primary school. That is why we 
present the development and validation of the competent CT test (cCTt), an unplugged CT test targeting 7-9 year-old students. In 
the first phase, 37 experts evaluated the validity of the cCTt through a survey and focus group. In the second phase, the test 
was administered to 1519 students. We employed Classical Test Theory, Item Response Theory, and Confirmatory Factor Analysis to 
assess the instruments' psychometric properties. The expert evaluation indicates that the cCTt shows good face, construct, and 
content validity. Furthermore, the psychometric analysis of the student data demonstrates adequate reliability, difficulty, and 
discriminability for the target age groups. Finally, shortened variants of the test are established through Confirmatory Factor 
Analysis. To conclude, the proposed cCTt is a valid and reliable instrument, for use by researchers and educators alike, which 
expands the portfolio of validated CT assessments across compulsory education. Future assessments looking at capturing CT in a 
more exhaustive manner might consider combining the cCTt with other forms of assessments.},
	number = {7},
	urldate = {2025-10-07},
	journal = {Journal of Educational Computing Research},
	author = {El-Hamamsy, Laila and Zapata-Cáceres, María and Barroso, Estefanía Martín and Mondada, Francesco and Zufferey, 
Jessica Dehler and Bruno, Barbara},
	month = dec,
	year = {2022},
	note = {arXiv:2203.05980 [cs]},
	keywords = {Computer Science - Computers and Society},
	pages = {1818--1866},
	annote = {Comment: To appear in the Journal of Educational Computing Research},
	file = {Preprint PDF:C\:\\Users\\cbenahmed\\Zotero\\storage\\B8L9FSD4\\El-Hamamsy et al. - 2022 - The competent 
Computational Thinking test (cCTt) Development and validation of an unplugged Computa.pdf:application/pdf},
}

@article{adorni_designing_2025,
	title = {Designing the virtual {CAT}: {A} digital tool for algorithmic thinking assessment in compulsory education},
	volume = {45},
	issn = {22128689},
	shorttitle = {Designing the virtual {CAT}},
	url = {http://arxiv.org/abs/2408.01263},
	doi = {10.1016/j.ijcci.2025.100760},
	abstract = {Algorithmic thinking (AT) is a critical skill in today's digital society, and it is indispensable not only 
in computer science-related fields but also in everyday problem-solving. As a foundational component of digital education and 
literacy, fostering AT skills is increasingly relevant for all students and should become a standard part of compulsory 
education. However, successfully integrating AT into formal education requires effective teaching strategies and robust and 
scalable assessment procedures. In this paper, we present the design and development process of the virtual Cross Array Task 
(CAT), a digital adaptation of an unplugged assessment activity aimed at evaluating algorithmic skills in Swiss compulsory 
education. The development process followed iterative design cycles, incorporating expert evaluations to refine the tool's 
usability, accessibility and functionality. A participatory design study played a dual role in shaping the platform. First, it 
gathered valuable insights from end users, including students and teachers, to ensure the tool's relevance and practicality in 
classroom settings. Second, it facilitated the collection and preliminary analysis of data related to students' AT skills, 
providing an initial evaluation of the tool's assessment capabilities across various developmental stages. This was achieved 
through a pilot study involving a diverse group of students aged 4 to 12, spanning preschool to lower secondary school levels. 
The resulting instrument features multilingual support and includes both gesture-based and visual block-based programming 
interfaces, making it accessible to a broad range of learners. Findings from the pilot study demonstrate the platform's 
usability and accessibility, as well as its suitability for assessing AT skills, with preliminary results showing its ability to 
cater to diverse age groups and educational contexts.},
	urldate = {2025-10-07},
	journal = {International Journal of Child-Computer Interaction},
	author = {Adorni, Giorgia and Piatti, Alberto},
	month = sep,
	year = {2025},
	note = {arXiv:2408.01263 [cs]},
	keywords = {Computer Science - Artificial Intelligence, Computer Science - Computers and Society, Computer Science - 
Human-Computer Interaction},
	pages = {100760},
	file = {Preprint PDF:C\:\\Users\\cbenahmed\\Zotero\\storage\\VXDPAJ4K\\Adorni et Piatti - 2025 - Designing the virtual 
CAT A digital tool for algorithmic thinking assessment in compulsory educatio.pdf:application/pdf},
}

@article{battal_computer_2021,
	title = {Computer {Science} {Unplugged}: {A} {Systematic} {Literature} {Review}},
	volume = {50},
	issn = {0047-2395, 1541-3810},
	shorttitle = {Computer {Science} {Unplugged}},
	url = {https://journals.sagepub.com/doi/10.1177/00472395211018801},
	doi = {10.1177/00472395211018801},
	abstract = {The computer science (CS) unplugged approach intends to teach CS concepts and computational thinking skills 
without employing any digital tools. The current study conducted a systematic literature review to analyze research studies that 
conducted investigations related to implementations of CS unplugged activities. A systematic review procedure was developed and 
applied to detect and subsequently review relevant research studies published from 2010 to 2019. It was found that 55 research 
studies (17 articles + 38 conference proceedings) satisfied the inclusion criteria for the analysis. These research studies were 
then examined with regard to their demographic characteristics, research methodologies, research results, and main findings. It 
was found that the unplugged approach was realized and utilized differently among researchers. The majority of the studies used 
the CS unplugged term when referring to “paper–pencil activities,” “problem solving,” “storytelling,” “games,” “tangible 
programming,” and even “robotics.”},
	language = {en},
	number = {1},
	urldate = {2025-10-07},
	journal = {Journal of Educational Technology Systems},
	author = {Battal, Ali and Afacan Adanır, Gülgün and Gülbahar, Yasemin},
	month = sep,
	year = {2021},
	pages = {24--47},
}

%Ref 1
@inproceedings{alayrangues_informatique_2017,
	title = {Informatique débranchée : construire sa pensée informatique sans ordinateur},
	shorttitle = {Informatique débranchée},
	abstract = {L'informatique débranchée, appelée également et peut-être plus justement, informatique sans ordinateur ou 
sciences manuelles du numérique, est une approche qui consiste à appréhender certains éléments de la science informatique par 
l'utilisation d'objets « concrets » et complètement « déconnectés » (bâtonnets, allumettes, cartes, jetons, ficelles, perle...). 
Elle permet de s'affranchir de la machine et de la technicité de sa programmation pour mieux saisir les grands principes de la 
science elle-même. Grâce à elle, il est possible d'aborder de manière ludique des problèmes complexes (recherche d'informations, 
tri de données, cryptographie, stratégie gagnante...) et leur résolution via la conception d'algorithmes. Elle favorise la mise 
en oeuvre d'une démarche scientifique avec les enfants et adolescents. Elle est propice au travail de groupe. Enfin, lors des 
activités, il s'agit non pas de donner une solution clé en main mais d'amener les participants à imaginer leur propre solution 
ou à la co-construire ensemble.},
	author = {S. Alayrangues, S. Peltier, and L. Signac},
	keywords = {ateliers d'initiation, informatique débranchée, médiation scientifique, science informatique},
	file = {HAL PDF Full Text:C\:\\Users\\cbenahmed\\Zotero\\storage\\66RJGDX4\\Alayrangues et al. - 2017 - Informatique 
débranchée  construire sa pensée informatique sans ordinateur.pdf:application/pdf},
}

%Ref 5
@article{sharif_university_of_technology_computational_2017,
	title = {Computational Thinking : A 21st Century Skill},
	volume = {11},
	shorttitle = {Computational Thinking},
	number = {2},
	author = {{Sharif University of Technology} and Tabesh, Yahya},
	pages = {65--70},
	file = {PDF:C\:\\Users\\cbenahmed\\Zotero\\storage\\78LUDTXV\\Sharif University of Technology et Tabesh - 2017 - 
Computational Thinking A 21st Century Skill.pdf:application/pdf},
}

